%!TEX program = lualatex

\documentclass[11pt]{article}

\usepackage[a4paper,margin=2.5cm]{geometry}
\usepackage{amsmath}
\usepackage{amssymb}

% Paragraphs without indent, separated by a little space instead
\setlength{\parindent}{0pt}
\setlength{\parskip}{5pt plus 1pt minus 1pt}

% Long code names cannot be hyphenated: stretch the spaces of their lines rather than overflow
\setlength{\emergencystretch}{3em}

\newcommand{\R}{\mathbb{R}}
\newcommand{\dt}{\Delta t}

\title{MonoSLAM}
\date{}

\begin{document}

\maketitle

MonoSLAM is a monocular EKF SLAM in the line of Davison et al.~\cite{davison2007}. It uses one
camera and no other sensor, and a single extended Kalman filter holds both the camera and the map.
A chessboard starts it and fixes the scale, and from then on it tracks the natural features of the
scene.

Section~\ref{sec:loop} describes the frame loop. Sections~\ref{sec:state} to~\ref{sec:idp} cover
the state, the motion model and the features. Sections~\ref{sec:flow} and~\ref{sec:search} explain
how the features are found in each frame, and section~\ref{sec:update} how the update uses them.
Section~\ref{sec:map} covers the map, and section~\ref{sec:global} how it is found again when
tracking fails. Section~\ref{sec:results} gives the results, and section~\ref{sec:limit} explains
why a filter cannot go much further.

\section{The frame loop}
\label{sec:loop}

MonoSLAM asks the camera for $320 \times 240$ frames, which it delivers at about 100~Hz with a
fixed exposure of 5~ms, and works on them without resizing (\texttt{FrameSource}). All the lengths
in pixels in this document refer to these frames, including the ones first measured on
$640 \times 480$ frames, which have been halved. A quarter of the pixels makes each frame about
twice as fast to process, with the same accuracy (section~\ref{sec:results}).

The camera is calibrated on $640 \times 480$ frames. This is the largest frame size the camera
gives with the same field of view, and the corners are more precise there. A larger mode with the same aspect ratio does not
necessarily have the same field of view: at $1024 \times 768$ this camera sees $12\%$ more in each
direction. The calibration is then rescaled to the small frames: a pixel center $x$ becomes
$(x + \frac{1}{2}) s - \frac{1}{2}$, where $s$ is the scale factor, and the distortion, which acts
on normalized coordinates, stays the same.

Every frame goes through the same steps, in \texttt{main.cpp}:
\begin{enumerate}
  \item \texttt{ekf.predict(dt)} moves the camera with the motion model over the real time between
        the two frames, taken from their capture times. Live, frames are dropped while one is being
        processed, so this interval can span several frame periods.
  \item \texttt{tracker.predict(gray)} uses the optical flow to move the \emph{seed} of every
        feature, the place where it is expected in the new image (section~\ref{sec:flow}).
  \item If the previous frame was weak (section~\ref{sec:global}), \texttt{keyframes.match}
        searches the features where the most similar keyframes saw them. If that fails,
        \texttt{matchGlobally} searches them anywhere in the frame by their appearance. When either
        succeeds, the camera is reset to the pose found. Each feature that was found gets its match
        as its new seed, and every other feature gets its projection from that pose.
  \item \texttt{ekf.observe()} predicts which features are in view and where. \texttt{matchFeatures}
        searches each of them around its seed, or around its prediction if it is coming back into
        view (section~\ref{sec:search}), and \texttt{ekf.update} decides with 1-point RANSAC which
        matches to use (section~\ref{sec:update}).
  \item If the update used at least four features, each seed is replaced by the projection of its
        feature from the updated state. Otherwise the seeds keep the positions given by the
        optical flow.
  \item Unless the frame is weak, the map is maintained. Features that are rarely matched are
        removed, and new ones are added where the frame has none (section~\ref{sec:map}). The frame
        also becomes a keyframe if it was taken far enough from all the others.
\end{enumerate}

Together, these steps form a loop: images, seeds, search regions, matches, filter, projections,
and back to the seeds. Step~5 closes the loop through the filter, which keeps the seeds from
drifting. Step~2 keeps the loop running through the optical flow when the filter cannot.

\begin{center}
\begin{tabular}{lll}
  step & reads & writes \\
  \hline
  \texttt{ekf.predict} & state & state \\
  \texttt{tracker.predict} & images, seeds & seeds \\
  \texttt{Keyframes::match}, \texttt{matchGlobally} & image, state, keyframes & state, seeds \\
  \texttt{matchFeatures} & image, predicted state, seeds & matches \\
  \texttt{ekf.update} & state, matches & state \\
  \texttt{tracker.correct} & state & seeds \\
  \texttt{Mapper} & image, state & state, seeds \\
  \texttt{Keyframes::add} & image, state & keyframes \\
\end{tabular}
\end{center}

All the information reaches the filter through its own models and its own update. The seeds only
determine where the features are searched. The relocalization of step~3 is the only place where the camera is set
without an update.

\section{State and conventions}
\label{sec:state}

The state is
\begin{equation}
  x = \begin{pmatrix} r \\ e \\ v \\ \omega \\ y_1 \\ \vdots \\ y_N \end{pmatrix} \in \R^{12 + 6N},
\end{equation}
where $r \in \R^3$ is the camera position in world coordinates, $e \in \R^3$ an orientation error
(see below), $v$ the linear velocity in world coordinates, $\omega$ the angular velocity in camera
coordinates, and $y_i \in \R^6$ the features, in inverse depth (section~\ref{sec:idp}). The world
frame is that of the chessboard, which lies on the plane $z = 0$. Lengths are in cm and angles in
rad.

\paragraph{The orientation.}
The orientation is a rotation matrix $R$, the change of basis from camera to world coordinates. A
point with camera coordinates $X_c$ has world coordinates
\begin{equation}
  X_w = R\, X_c + r,
\end{equation}
so the columns of $R$ are the camera axes expressed in the world, and $R_{cw} = R^\top$ goes from
world to camera. $R$ is not part of the state vector. The filter keeps it alongside, and the state
holds a small rotation $e$ of the camera on top of it, so that the orientation is
\begin{equation}
  R\, \exp([e]_\times),
\end{equation}
where $\exp([e]_\times)$ is the rotation by $\|e\|$ radians about $e$ (Rodrigues' formula,
\texttt{rotation}) and $[e]_\times$ is the cross-product matrix (\texttt{skew}). $P$ holds the
uncertainty of $e$, not of $R$. After every update, $R$ absorbs $e$,
$R \leftarrow R \exp([e]_\times)$, and $e$ is reset to zero. All the Jacobians are therefore taken
at $e = 0$, where $\exp([e]_\times) \approx I + [e]_\times$. This is the error-state, or
multiplicative, EKF~\cite{lefferts1982}.

When $R$ absorbs $e$, the covariance of the error must also be expressed about the new
reference. An
error $d$ about $R \exp([e]_\times)$ is an error $G\, d$ about $R$, with
$G = I - \frac{1}{2} [e]_\times$, so the orientation rows and columns of $P$ are multiplied by $G$.
For the small $e$ of a single update this looks negligible, but it is first order in $e$. Without
it, runs diverge within ten frames from those of a quaternion filter, which applies the same
correction implicitly.

Unlike a quaternion, $e$ has exactly as many entries as the orientation has degrees of freedom, so
there is no unit norm to maintain and $P$ is never singular. Unlike a rotation vector for the whole
orientation, $e$ stays small, far from the angle of $\pi$ where a rotation vector wraps around and
its covariance stops making sense.

\paragraph{Composing rotations.}
When the camera rotates there are two camera frames, the previous one $C$ and the new one $C'$.
The new orientation $R' = R \exp([w]_\times)$ gives
\begin{equation}
  X_w = R'\, X_{c'} + r = R\, \exp([w]_\times)\, X_{c'} + r,
\end{equation}
a chain of two changes of basis: $\exp([w]_\times)$ maps coordinates in $C'$ to coordinates in
$C$, and $R$ maps those to the world. The rotation leaves its axis $w$ fixed, so $w$ has the same
coordinates in $C$ and $C'$, and it is a rotation in camera coordinates. In the world it is
$R\, w$, and $R \exp([w]_\times) = \exp([R w]_\times) R$. The motion model composes rotations on
the right, so $\omega$ is in camera coordinates, and so is $e$.

\paragraph{The right Jacobian.}
To first order, a small change $d$ of the rotation vector of $\exp([\theta]_\times)$ is a small
rotation in its own frame,
\begin{equation}
  \exp([\theta + d]_\times) \approx \exp([\theta]_\times) \exp([J_r(\theta)\, d]_\times),
  \qquad
  J_r(\theta) = I - \frac{1 - \cos\varphi}{\varphi^2} [\theta]_\times
              + \frac{\varphi - \sin\varphi}{\varphi^3} [\theta]_\times^2,
\end{equation}
with $\varphi = \|\theta\|$ (\texttt{rightJacobian}). $J_r(0) = I$, and for the rotation during
one frame, $\varphi \approx 0.05$, it differs from $I$ by a few percent.

\section{The motion model}
\label{sec:motion}

Between two frames the velocities change by an unknown amount, the result of a random acceleration
that stays constant during the step: $V = a\, \dt$ and $\Omega = \alpha\, \dt$, with
$a \sim \mathcal{N}(0, \sigma_a^2 I_3)$ and $\alpha \sim \mathcal{N}(0, \sigma_\alpha^2 I_3)$. The
velocities also decay towards zero, with time constants $\tau_v$ and $\tau_\omega$. With
$e_v = e^{-\dt/\tau_v}$ and $h_v = \tau_v (1 - e_v)$, and likewise $e_\omega$ and $h_\omega$,
\begin{equation}
  \begin{aligned}
    r_{k+1} &= r_k + h_v\, v_k + \tfrac{1}{2} V\, \dt, \\
    R_{k+1} &= R_k \exp\big(\big[h_\omega\, \omega_k + \tfrac{1}{2} \Omega\, \dt\big]_\times\big), \\
    v_{k+1} &= e_v\, v_k + V, \\
    \omega_{k+1} &= e_\omega\, \omega_k + \Omega, \\
    y_{i,k+1} &= y_{i,k},
  \end{aligned}
  \qquad
  Q = \dt^2
  \begin{pmatrix} \sigma_a^2 I_3 & 0 \\ 0 & \sigma_\alpha^2 I_3 \end{pmatrix},
  \label{eq:motion}
\end{equation}
where $Q$ is the covariance of $(V, \Omega)$ (\texttt{MotionModel}).

\paragraph{The decay.}
$h$ is the distance a decaying velocity covers in one step, divided by its value at the start of
the step. As $\tau \to \infty$, $e \to 1$ and $h \to \dt$, which is the constant velocity model of
Davison et al.~\cite{davison2007}. As $\tau \to 0$ the velocity vanishes immediately, and the
camera is predicted to stay where it is. A hand-held camera keeps moving between consecutive
frames, but its velocity does not stay the same for long. The decay also limits how far the prediction can drift
when nothing is measured.

\paragraph{The factor $\frac{1}{2}$.}
If the acceleration is constant during the step, the velocity changes linearly from $v_k$ to
$v_k + V$, and the position moves by the mean velocity over the step,
$r_{k+1} = r_k + (v_k + \frac{1}{2} V)\, \dt$. The model of Davison et al.~\cite{davison2007} uses
$(v_k + V)\, \dt$ instead, an impulse at the start of the step. Both predict the same mean. They
differ in how much of the noise reaches the position: a variance of $\frac{1}{4} \sigma_a^2 \dt^4$
here, against $\sigma_a^2 \dt^4$. The rotation gets the same factor, which is exact for a fixed
axis and correct to first order otherwise.

\paragraph{Linearization.}
The prediction is linearized at the mean angular velocity, which is not zero. With
$\theta = h_\omega\, \omega_k$, the orientation becomes $R_{k+1} = R_k \exp([\theta]_\times)$, and
the error after the step is the $e_{k+1}$ that satisfies
\begin{equation}
  R_{k+1} \exp([e_{k+1}]_\times) = R_k \exp([e_k]_\times) \exp([\theta + h_\omega\, d\omega]_\times),
\end{equation}
for an error $d\omega$ of the angular velocity. Using the right Jacobian, and moving
$\exp([e_k]_\times)$ across $\exp([\theta]_\times)$ as above, gives
$e_{k+1} \approx \exp(-[\theta]_\times)\, e_k + h_\omega J_r(\theta)\, d\omega$. The Jacobian $F$
of the camera part is therefore the identity except for
\begin{equation}
  \frac{\partial r_{k+1}}{\partial v_k} = h_v I_3,
  \quad
  \frac{\partial e_{k+1}}{\partial e_k} = \exp([\theta]_\times)^\top,
  \quad
  \frac{\partial e_{k+1}}{\partial \omega_k} = h_\omega J_r(\theta),
  \quad
  \frac{\partial v_{k+1}}{\partial v_k} = e_v I_3,
  \quad
  \frac{\partial \omega_{k+1}}{\partial \omega_k} = e_\omega I_3.
\end{equation}
The noise enters through the velocities, so $W = \partial f / \partial (V, \Omega)$ is
$12 \times 6$, with $\frac{1}{2} \dt\, I_3$ and $I_3$ in the rows of $r$ and $v$, and
$\frac{1}{2} \dt\, J_r(\theta)$ and $I_3$ in those of $e$ and $\omega$. Only the camera moves, so
$P$ becomes
\begin{equation}
  \begin{pmatrix} F P_{cc} F^\top + W Q W^\top & F P_{cf} \\ P_{fc} F^\top & P_{ff} \end{pmatrix},
\end{equation}
where $c$ stands for the camera and $f$ for the features. $P$ is then symmetrized, so that rounding
errors do not accumulate over thousands of steps. The tests \emph{Process Jacobian} and
\emph{Process noise Jacobian} compare $F$ and the orientation block of $W$ with finite differences
of the prediction, with the error measured from the nominal prediction.

\paragraph{Values.}
$\sigma_a = 1000$~cm/s$^2$, $\sigma_\alpha = 15$~rad/s$^2$, $\tau_v = 0.3$~s and
$\tau_\omega = 0.2$~s. In one frame the velocities can change by about $33$~cm/s and $0.5$~rad/s,
and without measurements they halve in about $0.2$ and $0.14$~s. The time between frames comes from
their capture times, or is $1/30$~s for recordings that have none.

\paragraph{Why this model.}
The chessboard ground truth (section~\ref{sec:results}) gives the camera pose in every frame where
the whole board is detected. Running a Kalman filter on these poses with each candidate model, with
the noise fitted by maximum likelihood over three sequences, measures how well the model predicts
the next frame (negative log-likelihood, lower is better):

\begin{center}
\begin{tabular}{lrr}
  model & position & orientation \\
  \hline
  no velocity, a random displacement & $1827$ & $-10565$ \\
  impulse~\cite{davison2007} & $674$ & $-11655$ \\
  constant acceleration during the step & $649$ & $-11710$ \\
  same, with decay, $\tau = 10$~s & $639$ & $-11718$ \\
  same, with decay, $\tau = 0.3$~s & $595$ & $-11858$ \\
\end{tabular}
\end{center}

Having velocities in the state matters most, by more than a thousand: the motion of the camera
persists from one frame to the next. How the noise enters them matters little, and a decay of about
$0.3$~s fits best, at the cost of one more parameter. The motion itself only needs about
$90$~cm/s$^2$ and $5.5$~rad/s$^2$, but the filter needs larger values, for two reasons. First, the
noise also has to absorb the errors of the filter. Second, the fastest frames are the most blurred,
so the board is not detected in them and this fit never sees them.

\section{Inverse depth features}
\label{sec:idp}

A feature uses the inverse depth parametrization of Civera et al.~\cite{civera2008}:
\begin{equation}
  y = (a, \theta, \phi, \rho) \in \R^6,
  \qquad
  p = a + \frac{1}{\rho}\, m(\theta, \phi),
  \qquad
  m(\theta, \phi) = \begin{pmatrix} \cos\phi \sin\theta \\ -\sin\phi \\ \cos\phi \cos\theta \end{pmatrix},
\end{equation}
where the anchor $a$ is the position of the camera when the feature was first seen, $\theta$ and
$\phi$ are the azimuth and elevation of the ray from $a$ to the point, in world coordinates, and
$\rho$ is the inverse of the distance along that ray.

\paragraph{Why.}
A point seen once is only a ray: its depth stays unknown until the camera has moved enough. Along
the ray the uncertainty of the depth is very asymmetric, since a point at $50$~cm may be at $25$~cm
or at infinity. A Gaussian on $p$ cannot describe that, but a Gaussian on $\rho$ can:
$\rho \in [0, 0.04]$ covers everything from $25$~cm to infinity. The projection is also much
closer to linear in $\rho$ than in the depth. A feature can therefore enter the filter in the frame
where it is first seen, with a large $\sigma_\rho$, and constrain the orientation of the camera
long before its depth is known.

\paragraph{Observation.}
The vector
\begin{equation}
  h_c = R^\top \big( \rho\, (a - r) + m(\theta, \phi) \big)
\end{equation}
is the point in camera coordinates, scaled by $\rho$. The scale does not change its pixel, and it
keeps $h_c$ well defined at $\rho = 0$, for points at infinity. The pixel is the projection of $h_c$
through the lens distortion of the camera, and its Jacobians are those of $h_c$ with respect to
$r$, $e$ and $y$, multiplied by the Jacobian of the projection. Rotating the camera by $e$ turns
$h_c$ into $\exp(-[e]_\times) h_c \approx h_c + [h_c]_\times e$, so
$\partial h_c / \partial e = [h_c]_\times$ (\texttt{inverse\_depth::toCamera},
\texttt{CameraModel::projectCameraPoint}). A feature is in view if $h_c$ is in front of the camera
and its pixel falls inside the image.

\paragraph{No conversion to points.}
Civera et al.~\cite{civera2008} convert a feature to a point $p \in \R^3$ once its depth is known
well enough, which saves three entries of the state. Here features are never converted, and
keep their six parameters until they are removed. The observation model, the map and the code stay uniform, at the
price of a slightly larger state. \texttt{EKF::isWellLocated} identifies the features whose depth
is known, $\rho > 0$ and $\sigma_\rho < 0.1\, \rho$, which is where Civera et al.\ would convert
them. \texttt{inverse\_depth::toPoint} gives the point when one is needed, for drawing or for
warping patches. A feature with $\rho \le 10^{-3}$ is placed $10$~m along its ray.

\paragraph{Initialization.}
The first camera pose comes from the whole chessboard, with PnP, and is accepted if it explains all
the corners within $3$~px and with an RMS below $1$~px. Five of the corners become the first
features: the four outermost ones and the one closest to the center. They are placed at their
known positions $p_i$, with the anchor at the initial camera position,
$y_i = (r_0, \theta_i, \phi_i, \rho_i)$ (\texttt{fromPoint}). The covariance of
$(\theta_i, \phi_i, \rho_i)$ is $J\, (0.05\, I_3)\, J^\top$, where $J$ is the Jacobian from
\texttt{fromPointJacobian}, which amounts to a variance of $0.05$~cm$^2$ on each axis of the point.
The anchor has no uncertainty, since it is a constant and not a measurement. The camera starts with
standard deviations of $0.22$~cm, $0.02$~rad, $30$~cm/s and $1.5$~rad/s, and no velocity. This is
where the scale of the map comes from. The chessboard plays no further role.

\paragraph{New features.}
A feature first seen at pixel $u$ from the pose $(r, R)$, with an inverse depth $\rho_0$, is
\begin{equation}
  y = \big(r,\; \theta(d),\; \phi(d),\; \rho_0\big), \qquad d = R\, \pi^{-1}(u),
\end{equation}
where $\pi^{-1}(u)$ is the ray of the pixel in camera coordinates, found by undistorting it.
Rotating the camera by $e$ turns $d$ into $R \exp([e]_\times) \pi^{-1}(u) \approx d - R\,
[\pi^{-1}(u)]_\times e$. With $J = \partial y / \partial (r, e, u, \rho_0)$, of size $6 \times 9$
(\texttt{inverse\_depth::fromPixel}), the state grows to
\begin{equation}
  x \leftarrow \begin{pmatrix} x \\ y \end{pmatrix},
  \qquad
  P \leftarrow
  \begin{pmatrix}
    P & P_{:, c}\, J_c^\top \\
    J_c\, P_{c, :} & J_c\, P_{cc}\, J_c^\top + J_u\, \sigma_u^2\, J_u^\top + J_\rho\, \sigma_\rho^2\, J_\rho^\top
  \end{pmatrix},
\end{equation}
where $c$ stands for the six entries of the camera position and orientation error, $\sigma_u$ is the
pixel noise, and $\sigma_\rho$ the uncertainty of the inverse depth (\texttt{EKF::addFeature}).
Because the new feature is correlated with the pose it was added from, a later update that
corrects that pose also corrects the feature.

\section{The seeds}
\label{sec:flow}

Every feature has to be found again in each new image, and the first question is where to look. In
EKF SLAM the search region is the $3\sigma$ ellipse of the innovation covariance around the
prediction, the active search of Davison et al.~\cite{davison2007}. That only works if the ellipse
contains the feature and nothing that looks like it, and a hand-held camera jerks. On the first
sequence the chessboard corners move up to $19$~px between frames, and a constant velocity
prediction only brings the median error down from $5$ to $3$~px, with $11$~px at the $99$th
percentile. With the active search, the first two sequences are lost in the frames where the motion
changes by $10$~px from one frame to the next. The motion model cannot predict such changes, but
the images can.

\texttt{FlowTracker} follows the image motion of every feature, its \emph{seed}, without any
knowledge of the filter. It keeps the previous image and one seed per feature, including those out
of view, and in every frame (\texttt{predict}) it does three things:
\begin{enumerate}
  \item \emph{It tracks generic corners.} The corners of the previous image are tracked into the
        new image and back with pyramidal Lucas--Kanade~\cite{lucas1981}, with six levels and an
        $11 \times 11$ window. A corner is kept if it comes back within $0.5$~px of where it
        started, which removes most corners lost to blur or occlusion. The corners carry over from
        frame to frame, and new ones are only detected when fewer than $120$ are left
        (\texttt{goodFeaturesToTrack}, quality $0.01$, at least $4$~px from each other and from the
        remaining ones), up to $150$. Detecting them again in every frame made the flow $40\%$
        slower. The six levels matter: with four, the second sequence is lost, because the flow
        cannot follow jumps of more than $50$~px.
  \item \emph{It moves each seed with its neighbors.} An affine map is fitted with RANSAC to the
        motion of the $24$ tracked corners closest to the seed (inliers within $1.5$~px), and the
        seed is moved by that map. With fewer than $6$ corners, or if the fit fails, the seed stays where
        it is. Seeds outside the image are moved in the same way, by the corners closest to them.
  \item \emph{It lets the seed's own flow take over.} The seeds in view are also tracked
        themselves, with a $15 \times 15$ window. If the result lands within $3$~px of the affine
        position, it replaces it, since the flow of a single point is more precise. If not, the
        seed has probably slipped to a similar-looking neighbor, and the affine position is kept.
\end{enumerate}
An affine map is the first-order approximation of the image motion around a point. Using it assumes
that the seed moves like the corners around it, which is true for a locally smooth surface, and for
any scene while the camera mostly rotates. At a depth discontinuity, with the camera translating,
the fit follows whichever surface has more corners. The test \emph{Flow follows depth layers}
checks that two regions moving apart each carry their own seeds.

After the update, if at least four features were used, \texttt{correct} replaces each seed with
the projection of its feature from the updated state. A seed that slipped to a neighbor is
corrected this way, since the other features fix the pose. With fewer features the pose is not
well constrained, and the seeds keep the positions given by the optical flow.

\section{The patch search}
\label{sec:search}

Each feature keeps the $11 \times 11$~px patch around it in the frame where it was first seen. To
find it in a new frame (\texttt{findPatch}), the patch is first warped to how it would look from
the predicted pose. The patch is treated as a piece of the plane through the feature with its
normal: the chessboard's for the five chessboard corners, and facing the first camera for all other
features. The corners of the patch are mapped through the whole camera model, distortion included,
and the warp is the homography that maps them.

The warped patch is searched within $\pm 8$~px of the seed, or of the prediction if the seed is out
of view, by zero-mean normalized cross-correlation over the mask of the warped patch. The best peak
is kept if it scores at least $0.6$ and no other peak, meaning a local maximum more than $1$~px
away, scores within $0.05$ of it. Two nearly equal peaks mean a repeated texture or an edge, and
which one wins then depends on tiny differences in the prediction. Every feature found is returned,
and the update decides which ones to use.

\paragraph{Subpixel alignment.}
A parabola through the peak and its neighbors gives a first subpixel position, but it is biased
towards whole pixels. The warped patch is then aligned with the frame by the inverse compositional
Lucas--Kanade algorithm~\cite{baker2004} for a translation, with at most $10$ iterations
(\texttt{alignPatch}). At every step the frame is brought to the mean and contrast of the patch, as
the correlation does. The alignment is abandoned if it wanders more than $1$~px away from the peak.
It succeeds for $84$ to $99$\% of the matches and moves them by a median of $0.12$ to $0.18$~px.
The fit also tells how well the patch was aligned. Its covariance is the variance of the residuals
over the Hessian of the patch gradients,
\begin{equation}
  \Sigma_a = \frac{\sum_k e_k^2}{n - 2} \Big(\sum_k \nabla T_k \nabla T_k^\top\Big)^{-1},
\end{equation}
which is large when the patch does not look like the frame, or when it lacks texture in some
direction, as along an edge. For the median match, its larger standard deviation is $0.1$~px. On
its own, the alignment does not change the results beyond the spread of the runs: the precision of
the matches is not what limits them. Its covariance does improve them (section~\ref{sec:update}).

\paragraph{Returning features.}
A feature that was out of view for at least $10$ frames is \emph{returning} until the update uses
it again. Its seed is useless, because the flow could not follow the feature while it was away, so
it is searched around its prediction instead. The region is $3\sigma$ of the larger standard
deviation of its innovation, as in the active search of Davison et al.~\cite{davison2007}, but no
smaller than $\pm 8$~px and no larger than $\pm 20$~px. While the camera is away, the uncertainty of
the feature grows, and so does its search region, so the feature is found even if the map has
drifted. The upper limit keeps a large uncertainty from letting look-alikes into the region: without it, at
$\pm 30$~px, some runs of sequence~6 end more than $3$~m off. Over $15$ runs, searching returning
features this way recovers $17$ of the $30$ kidnaps, against $9$ when they were searched around
their seeds.

\section{The update}
\label{sec:update}

\texttt{EKF::observe} predicts the pixels $h$ of the features in view, two rows per feature. It
also computes their Jacobian $H$, the innovation covariance $S = H P H^\top + \sigma^2 I$ with
$\sigma = 1.5$~px, and the columns of $P$ that $H$ acts on. $H$ only has the columns of the camera
and of the features in view. A pixel depends only on the camera and on its own feature, so most of
$H$ is zero, and the product $P H^\top$ skips those zeros.

\paragraph{The noise of each match.}
$\sigma$ covers many things the alignment does not see, such as the error of the warp, and it
remains the noise of a well aligned match. A match aligned worse gets more noise, in the directions
where it was worse (\texttt{matchNoise}). With $\Sigma_a = V \operatorname{diag}(\lambda_1,
\lambda_2) V^\top$,
\begin{equation}
  R = \sigma^2 V \operatorname{diag}\big(\max(1, \lambda_1 / \lambda_0), \max(1, \lambda_2 / \lambda_0)\big) V^\top,
  \qquad \lambda_0 = 0.04~\text{px}^2,
\end{equation}
which replaces the block $\sigma^2 I$ of the match in $S$. About one match in ten has an alignment
standard deviation above $\sqrt{\lambda_0} = 0.2$~px. No match is trusted more than before, so this
can never make the filter more confident. Over $15$ runs the kidnaps recovered go from $18$ to $22$
of $30$, the returns of sequence~5 are as good or better, and everything else stays within the
spread of the runs. With $\lambda_0 = 0.01$~px$^2$ the kidnaps recovered were $23$, but the second
visit of sequence~6 was worse. With $0.0025$~px$^2$, where most matches get extra noise, they were
$15$, and sequence~3 was lost once.

\paragraph{One EKF step.}
\texttt{correct} applies the standard EKF update with a subset of the rows of the observation,
\begin{equation}
  x = x_0 + K \big(z - h(x_0)\big),
  \qquad
  K = P H^\top S^{-1},
  \qquad
  P \leftarrow P - K H P,
\end{equation}
solving $S K^\top = (P H^\top)^\top$ for $K^\top$ by Cholesky. Nothing is updated if $S$ is not
positive definite. $R$ then absorbs the orientation error (section~\ref{sec:state}).

\paragraph{1-point RANSAC.}
Some of the matches of a frame are wrong, and \texttt{EKF::update} decides which ones to use,
as in Civera et al.~\cite{civera2010}:
\begin{enumerate}
  \item \emph{Hypotheses.} For every match $j$, a copy of the state is updated with that match
        alone,
        \begin{equation}
          x^{(j)} = x_0 + K_j \big( z_j - h_j(x_0) \big),
          \qquad
          K_j = P H_j^\top S_{jj}^{-1}.
        \end{equation}
        A single pixel only fixes two degrees of freedom of the pose. The others move according to
        their correlations with those two in $P$, which is the most likely way given the
        prediction, and that is enough to tell which other matches agree. The hypotheses are states
        with $e \neq 0$, projected with $R \exp([e]_\times)$.
  \item \emph{Support.} The prediction $h(x^{(j)})$ is computed with the full nonlinear model, and
        every match within $3$~px of it supports the hypothesis. The features in view at $x^{(j)}$
        need not be the same as those in view at $x_0$. Rows are therefore paired by the feature
        they belong to, not by their position, and a match whose feature leaves the image does not
        support the hypothesis. The hypothesis with the largest support wins, and a tie is broken by
        the smallest sum of distances.
  \item \emph{Low innovation inliers.} The state is updated with the winning support.
  \item \emph{High innovation inliers.} A match can be right and still disagree with every
        single-match hypothesis, if the prediction was far off. After the first update the
        prediction is much better, so $h$, $H$ and $S$ are computed again. Every remaining match
        with a normalized innovation squared below $9.21$, the $99\%$ quantile of $\chi^2_2$, is
        rescued and used in a second update. Updating in two steps is equivalent to updating with
        both sets at once, up to the linearization.
\end{enumerate}
The thresholds, $3$~px and $9.21$, are part of the measurement model, together with $\sigma$
(\texttt{MeasurementModel}). Since everything goes through $h$, $H$, $P$ and $\sigma$, nothing in
the update knows about poses or images. Each hypothesis costs one evaluation of the observation
model, which is cheap for a dozen features. With hundreds, a random subset of hypotheses would do,
as in Civera et al.~\cite{civera2010}.

\paragraph{The plain update is enough.}
The projection is not linear, and an iterated EKF, which is Gauss--Newton on the same cost, would
follow its curvature. Without a line search it oscillates on this cost. With the few rescued
matches of the second stage, which do not determine a pose on their own, it wanders along the
directions they leave free. And the curvature it would correct is small for the innovations that
1-point RANSAC lets through.

\section{The map}
\label{sec:map}

\texttt{Mapper} maintains the map at the end of every frame that is not weak.

\paragraph{Adding.}
While the update of the frame used fewer than $12$ of the features in view, up to two new ones are
added per frame. A feature predicted in view but not found does not count:
\begin{enumerate}
  \item \emph{Where.} A box a quarter of the frame in each dimension slides over it in steps of half
        its size, staying $10$~px from the border. The positions with no feature in view inside are
        tried in order, the farthest from the existing features first, so that the features keep a
        roughly constant density over the frame. Davison et al.~\cite{davison2007} place the box at
        random instead, but a fixed order keeps the runs repeatable. A box can be empty because
        there is nothing to track in it, and then the next one is tried.
  \item \emph{Salient.} \texttt{goodFeaturesToTrack} proposes up to $10$ corners in the box
        (Shi--Tomasi~\cite{shi1994}, quality $0.01$, at least $5$~px apart), $15$~px away from the
        features in view, strongest first.
  \item \emph{Textured.} The quality of \texttt{goodFeaturesToTrack} is relative to the best corner
        in the box, so on a blank wall the least blank point still passes. A candidate also needs
        an absolute amount of texture: the root mean square of the gradient of its patch along its
        least textured direction. This is the smaller eigenvalue of the gradient matrix of the
        patch, as in Shi and Tomasi, and it must be at least $4$ gray levels per pixel. On the six
        sequences, the features below $2$, $31\%$ of them, were kept for a median of $11$ frames and
        were used in $21\%$ of their searches, while those above $4$ were used in about $70\%$. The
        correlation used in the next test cannot detect a blank patch: it divides by the contrast
        of the patch, and what remains is camera noise.
  \item \emph{Distinctive.} A candidate is rejected if its patch correlates with some other place
        within $\pm 20$~px, outside $\pm 2$~px of itself, with a score of $0.8$ or more. This keeps
        new features off the inner corners of the chessboard, and off repeated texture in general,
        which the search could not tell apart within its region.
  \item \emph{Initial depth.} $\rho_0$ is the inverse of the median distance of the well located
        features in view, or $0.02$~cm$^{-1}$ ($50$~cm) if there are none. Its uncertainty is
        $\sigma_\rho = \rho_0$, so the $2\sigma$ interval covers everything from a third of that
        distance to infinity.
    \item The feature gets its patch, a descriptor (section~\ref{sec:global}) and a seed.
\end{enumerate}

\paragraph{Counting the features used.}
Counting all the features predicted in view lets the ones that fail keep the new ones out. When the
camera turns, the features about to leave the frame, blurred or seen from a new angle, are still
predicted in view: on sequences 7 and 8, in the frames where the update used fewer than $8$
features, a median of $13$ were predicted in view and only $7$ used. Over $15$ runs, counting the
features used recovers $26$ kidnaps of $30$ instead of $22$, closes the loop of sequence~6 at frame
$836$ in every run instead of $11$, and halves the error at the return of sequence~7, from $3.3$ to
$1.4$~cm, where it falls from $7.4$ to $3.6$ times the standard deviation the filter reports. It
costs a larger map, $15$ to $40\%$ more features at the end of a sequence, about $2$~ms per frame,
and one run of sequence~2 that is lost. Spreading the features further apart instead, at least
$30$, $45$ or $60$~px from those in view rather than $15$, did not help: with $45$ and $60$~px only
$14$ kidnaps of $30$ were recovered.

\paragraph{Removing.}
In every frame, each feature that is searched counts an attempt, and a match if the update used
it. After $10$ attempts, a feature used in less than half of its last $20$ is removed from the
state, the covariance (\texttt{EKF::removeFeatures}), the list of features and the tracker. Only
recent attempts count, so a feature that is no longer matched is removed within a third of a
second, even if it had been reliable for a long time. This removes occluded features, features on the edges of objects,
whose patch changes with the point of view, and features on things that move. A feature that comes
back into view after at least $10$ frames away starts counting again from zero. Its old failures
are no longer relevant, and a single bad frame on its return would otherwise remove it.

Two kinds of features are never removed: the chessboard corners, which hold the scale of the map,
and features used in at least $30$ frames. A feature used that often is part of the map. When
the camera returns after a long excursion, these features are what closes the loop. Their
first searches after the return may fail, because the state has drifted, and under the rule above
they would be gone before they could pull the state back. Without this exception, sequence~6 comes
back to the board after a minute with nothing left of the first visit except the board itself,
whose identical corners then match a whole square away.

\section{Finding the map again}
\label{sec:global}

The search only looks near the seeds and the predictions. When the camera gets somewhere else
without the images or the filter following it, as in a kidnap, simulated by cutting frames out of a
sequence, the features are no longer where the search looks for them. The features are then searched where the keyframes
that resemble the frame saw them, as in Klein and Murray~\cite{klein2008}
(\texttt{Keyframes::match}). If that fails, they are searched by their appearance alone,
anywhere in the frame, as in Williams et al.~\cite{williams2011} (\texttt{matchGlobally}).

\paragraph{Keyframes.}
After every frame that is not weak and used at least $8$ features, the frame becomes a keyframe if
no existing keyframe is within $5$~cm and $8^\circ$ of its pose. It is stored shrunk to
$80 \times 60$ pixels, blurred with $\sigma = 2.5$ of those pixels and with its mean removed,
together with the pose (\texttt{Keyframes::add}). After a weak frame, the frame is shrunk in the
same way, and the three keyframes with the smallest mean squared difference to it are tried in
turn. Since a small turn of the camera shifts the whole image, each keyframe is first aligned with the
frame by the shift that makes them most similar, within $6$ thumbnail pixels. The well located
features, except those of the chessboard, are projected with the pose of the keyframe and shifted
by the same amount. Each one is searched within $10$~px, with its patch warped to that pose. A pose is fitted to the features
found, and verified as described below for the descriptors.

The global match needs a detected corner to fall on the feature, and an ORB descriptor that stays
similar over the few pixels between the two. On these sequences only $3$ or $4$ of its $10$ to $14$ matches are
right, even while tracking is good, out of about $9$ features in view. The keyframes need neither,
because they search with the same patches that track the features, around where those should be.
They do need the camera to be back near a pose it had before. On sequence~5, right after the second
cut, the keyframe that looks most like the frame is $5$~cm and $4^\circ$ away from it. At the third
visit of sequence~6, the camera is $9$ to $14$~cm from every pose it had, and neither method finds
the map. Over $15$ runs, the kidnaps recovered went from $19$ to $23$ of $30$, and the return to the
board at frame $836$ of sequence~6 is within $10$~cm in $13$ runs instead of $9$.

\paragraph{Descriptors.}
Every feature gets an ORB descriptor~\cite{rublee2011} when it is created (\texttt{describe}). The
descriptor is made of $256$ binary comparisons of pixel pairs in the $31 \times 31$ patch around the
feature. The pairs are rotated by the orientation of the patch, the direction to the centroid of its
intensities, which turns with the image. ORB only computes this orientation for the corners it detects itself, so it is computed
here. Seen from the side, a feature can be $60$ bits away from how it looked when it was created.
So while a feature is tracked, its descriptor in the current frame is also kept if it is more than
$30$ bits away from all those it already has, up to eight (\texttt{Mapper::recordAppearance}).

\paragraph{Matching.}
Up to $1500$ corners of the frame are detected with \texttt{goodFeaturesToTrack}, as the features
were, since ORB's own corners rarely fall on them, and described in the same way. Each descriptor of
each well located feature is matched to its nearest corner. The match is kept if that corner is
within $64$ bits, and clearly closer than the best corner more than $5$~px away from it, by a ratio
of $0.8$. Corners
around the same structure look alike, but that is not an ambiguity. The five chessboard corners
are left out, since they all look the same. A pose is fitted to the matches with P3P and RANSAC,
with $1000$ iterations and inliers within $2$~px, and at least six matches must agree.

\paragraph{Verification.}
A wrong pose can agree with six descriptors by chance, so the pose is verified with the patches.
The features are projected with it and searched, warped to it, within $4$~px of their projections,
and at least six must be found within $1.5$~px.

\paragraph{When it is used.}
A frame is \emph{weak} when the update used fewer than four features, or fewer than half of those
searched. The state may then be wrong, even if some matches still agree with it: after a kidnap,
1-point RANSAC can keep finding a few look-alikes that agree with a wrong pose.
\begin{itemize}
  \item After a weak frame, the keyframes or the global match decide. If either finds a pose,
        \texttt{resetCamera} resets the camera to it, with no velocity, a large uncertainty
        ($2$~cm, $0.2$~rad, $30$~cm/s, $1.5$~rad/s) and no correlation with the features. The seeds
        are then replaced by the projections of their features.
  \item While fewer than half of the features searched are used, the map is left alone. No feature
        is added, since it would be placed from a wrong pose, and no failed match is counted, since
        the state, not the features, is to blame for the failures. A frame with fewer than four features used, but
        most of those searched, is different: the camera is entering a part of the scene that has
        no features yet, and it needs new ones. Freezing the map there would keep tracking weak
        and the map frozen indefinitely.
\end{itemize}

\section{Results}
\label{sec:results}

\paragraph{Ground truth.}
The sequences start on the chessboard, and some come back to it. \texttt{analyze\_pattern} finds
the whole board in every frame where it can, and fits the camera pose to its $54$ corners. A run is
scored in two ways. The first is in pixels: the distance from the detected corners to the
projections of the five chessboard features from the estimated pose
(\texttt{evaluate\_tracking.py}). The second is the error of the camera position when the board
comes back into view. The detector misses the blurriest frames,
which are the fastest ones, so the annotated frames are checked as well.

\paragraph{The runs are chaotic.}
A tiny difference, such as changing the pixel noise by one part in a billion, can flip a borderline
decision, such as whether a match is accepted or a feature added. From then on, the two runs diverge. A
single run cannot tell a better system from a lucky one, so every number below pools $15$ runs:
five sets of initial chessboard corners, each with three negligible changes of the pixel noise.

\begin{center}
\begin{tabular}{ll}
  sequence & result over $15$ runs \\
  \hline
    1 & median $0$ frames over $5$~px, at most $3$ \\
  2 & median $5$ frames over $5$~px, lost in $1$ run \\
  3 & median $3$ frames over $5$~px, at most $53$ \\
  5, back on the board at frames 336, 514, 566 & median errors $1.6$, $1.5$, $2.7$~cm \\
  5, kidnaps: frames 141--249 and 421--479 cut & recovered in $12$ and $14$ runs \\
  6, board visits at frames 359, 411 & median errors $4.4$, $1.8$~cm \\
  6, board visits at frames 771, 836 & median errors $3.0$, $1.7$~cm \\
  6, board visits at frames 1413, 1549 & lost in every run \\
\end{tabular}
\end{center}

Sequences 1 to 3 look at the board most of the time. Sequences 4 and 5 leave it after the
initialization and track natural features, such as a table, a bowl, a cable, chairs and pictures,
with $5$ to $20$ features per frame. Sequence~5 comes back to the board twice. Sequence~6 is a
minute long and comes back three times. The map holds the first time. The second excursion drifts,
often by tens of cm by the time the board comes back into view at frame 771. The features kept
from the first visit then close the loop: by frame 836, every run is within $10$~cm. Before
the third visit the camera sweeps a blurred, almost textureless sofa, and tracking is lost in every
run. The global match does not find the map again before the end. In the run lost on sequence~2, a
chessboard feature locks onto the wrong corner at frame $406$, and the five features then project
about $110$~px away from the board.

\paragraph{Speed.}
On sequence~6, while tracking, a frame took $6$~ms, the median of the averages over every second,
against $14$~ms with the same code on $640 \times 480$ frames. Since new features are added while
fewer than $12$ are used (section~\ref{sec:map}), the map is larger, and a frame takes $9$~ms on
this machine, against $7$~ms measured again without that change. At that size the filter took $1$ to
$3$~ms, and most of the time went to the images. The optical flow took $6$ to $10$~ms, of which
$4$~ms went to the Lucas--Kanade tracking of the $150$ corners forwards and backwards, and up to
$3$~ms to the affine motion fitted around each seed. The search of the patches took $2$ to $6$~ms.
With a quarter of the pixels, and windows and patches half as wide, the images cost a fraction of
that. Over $15$ runs, the accuracy is the same within the spread of the runs: the first return of
sequence~5 is $1$~cm worse, the second visit of sequence~6 is better, and the kidnaps recovered go
from $20$ to $18$ of $30$. The subpixel alignment of the matches adds about $1$~ms. Three more
things keep the time down:
\begin{itemize}
  \item The masked correlation of the patches is computed from three plain correlations
        (\texttt{maskedCorrelation}). It gives the same score as OpenCV's masked
        \texttt{matchTemplate}, $2.7$ times faster.
  \item The pyramid of each frame for the flow of the corners is built once, and used forwards,
        backwards, and again as the previous frame. It is built with the window of the corners. With the larger window of the features it would
        have one level less, and the corners would
        lose track of fast motions, as they did in $11$ of the $15$ runs of sequence~3.
  \item When new features are added, each candidate corner is checked only once. The boxes overlap
        and are searched again for every feature added, and checking every candidate each time took
        up to $40$~ms in one frame out of $14$.
\end{itemize}
After a weak frame, the keyframes and the global match run, so while the camera is lost a frame
takes longer, up to $30$~ms on average over a second. Live, a frame that takes longer than the
frame period causes the next one to be dropped, which doubles the interval of the prediction, and
with it the uncertainty of the predictions. With \texttt{--timing} the decoding of the camera
frames is reported too, about $1$~ms. Each frame that is read is a new image that nothing writes to
afterwards, so tracking, the features and the window share it without copying it. The window is
drawn and shown by the main thread, while tracking runs in another thread, which only hands over
its latest frame (\texttt{Display}). At $640 \times 480$ the display took $3$~ms per frame, and much
more with the window maximized, which was enough to drop frames. $150$ corners track as well as
$300$, but on $640 \times 480$ frames a $15 \times 15$ window with at most $10$ iterations lost
sequences 2 and 3 in every run.

\paragraph{Limitations.}
\begin{itemize}
  \item At a depth discontinuity, the seed of a feature follows the surface with more corners.
        Keeping both predictions, the flow of the feature itself and the motion of its neighbors,
        and letting the update choose, would fix it.
  \item $4$ of the $30$ kidnaps are not recovered in time, mostly the first one, where the camera
        sees, after the cut, a part of the scene it had not seen before. Which kidnaps are recovered
        changes a lot from run to run: the floating point rounding of the faster correlation alone
        moved the first kidnap from $13$ recoveries to $9$, and the second from $10$ to $11$, out
        of $15$.
  \item Over long excursions the map drifts, and the filter does not know it: the uncertainty it
        reports grows far less than the error. This is the limit of the filter itself
        (section~\ref{sec:limit}).
  \item The camera has a rolling shutter. On the chessboard frames, the reprojection error is
        smallest if the rows are exposed over about $10$~ms, from top to bottom. Predicting each
        feature at the time of its row did not help over $15$ runs, since the motion in these
        sequences is slow, and it only moves the corners by about $0.01$~px on average. It may
        matter for fast live motion. \texttt{evaluate\_tracking.py} would then need it as well,
        since it projects with a single pose per frame.
  \item The features used for a long time are never removed, so the map, and the cost of the
        filter, grow with the explored scene.
  \item The global match only describes detected corners, which lie a few pixels away from where
        the descriptor of a feature would match best, and ORB descriptors change a lot over a few
        pixels. On the return to the board at the end of sequence~6, most visible features match
        their descriptors well within $4$~pixels of their true position, but only one or two at the
        corners that are described.
\end{itemize}

\section{Where the filter ends}
\label{sec:limit}

When the camera turns away from the mapped part of the scene, every feature in view is new, and
the filter becomes confident about a pose that is wrong. This section shows where that comes from,
that none of the remedies from the filtering literature removes it, and that the images themselves
contain the information the filter loses.

\paragraph{A simulation.}
In a simulation, the models of the filter can be made exactly true. The camera follows the motion
model of section~\ref{sec:motion}, with its noise drawn as the filter assumes, among $20\,000$
points on a shell $1.5$ to $6$~m around the start. The pixels are the true projections plus
Gaussian noise with the $\sigma$ of the filter, data association is perfect, and features are added
and dropped as the mapper does, with the five chessboard corners known. The table gives the
normalized estimation error squared of the camera pose (NEES, which is $6$ for a consistent
filter), at $100$~Hz, over $20$ runs of $10$~s:

\begin{center}
\begin{tabular}{lrr}
  new features enter with & NEES, last $2$~s & position error \\
  \hline
  no measurements, prediction only & $5.1$--$5.5$ & \\
  their true inverse depth, $\sigma_\rho \approx 0$ & $5$--$6$ & $1.1$~cm \\
  an inverse depth drawn with $\sigma_\rho = \rho$, as in MonoSLAM & $721$ & $13.7$~cm \\
  the same, Jacobians at the true features & $5$--$7$ & $1.4$~cm \\
  an inverse depth drawn with $\sigma_\rho = 0.1\, \rho$ & $5$--$6$ & $1.6$~cm \\
\end{tabular}
\end{center}

With exact models and exact matches, the filter is consistent if the features enter with a known
depth. If they enter as MonoSLAM adds them, it is over a hundred times too confident, and reports
a position uncertainty of $1.4$~cm for an error of $13.7$~cm. Linearizing the observations of the
features at their true parameters, with the camera still at its estimate, makes it consistent
again. The cause is therefore the linearization of the young features at a wrong inverse depth. It
requires the camera to rotate and translate at the same time, and it happens with noiseless pixels
as well. Every update with a young feature gives the camera information computed with a wrong
Jacobian, and later updates cannot take it back.

\paragraph{The remedies of filtering.}
None of these brings the NEES down:
\begin{itemize}
  \item \emph{The second-order term} of the projection in $\rho$ and in the camera position, added
        to $S$, lowers it from $880$ to $620$ at $30$~Hz, and the position error stays the same.
  \item \emph{A partial (Schmidt) update}, in which the features with $\sigma_\rho > 0.3\, \rho$ only
        update themselves, with the Joseph form for the covariance, makes it worse. The update
        of the young feature itself still uses the Jacobian at its wrong $\rho$.
  \item \emph{First-estimate Jacobians} for the features make it worse, because the first estimate
        is the uncertain one.
  \item \emph{An iterated EKF}, relinearizing about each new estimate, is as consistent as the plain
        update when the depths are exact, and worse when they are not: $1460$ with $3$ iterations,
        and $2874$ with $10$. A single frame cannot observe $\rho$, so iterating only commits the
        camera harder to the wrong value. The linearizations that are wrong belong to the past
        frames.
  \item \emph{Delayed initialization} tracks the features outside the filter with Lucas--Kanade,
        fits their inverse depth along the epipolar line, and lets them in once
        $\sigma_\rho \le 0.15\, \rho$, or after $30$ frames. On the real sequences it lost to the
        current filter in every variant. When the camera enters a new part of the scene, the filter
        is left with $2$ to $5$ features instead of $8$ to $18$, and it drifts while it waits.
\end{itemize}
The invariant EKF fixes the inconsistency that comes from the unobservable directions of SLAM, the
global pose and the scale. Here the five chessboard features, which are never removed and have a
tiny covariance, already pin those down, and the Jacobians of an invariant EKF would still depend
on the depth of the young features.

\paragraph{Two sequences at $100$~Hz.}
Sequences 7 and 8 last $18$~s each. They start on the board, explore the room and come back. The
board is detected less reliably near the border of the frame, where its corners are fitted with an
RMS of $2$ to $4$~px instead of less than $1$, so only frames fitted within $1.5$~px are scored. In
the runs that close the loop, the position error is $1$ to $4$~cm, about $2$ to $3$ times the
standard deviation the filter reports. On sequence~7, two of the five runs, which differ only in
their initial corners, end $80$ to $120$~cm off. These runs, and the analysis below, predate the
count of the features used (section~\ref{sec:map}), which halves the error at the return of
sequence~7.

These two runs go wrong between frames $860$ and $1000$, when the board leaves the view and the
camera pans from the table to a bookshelf. By then all the features in view are new, and the five
runs do not agree on the motion:

\begin{center}
\begin{tabular}{lrr}
  run & displacement $860 \to 1000$, cm & bundle adjustment, cm \\
  \hline
  default corners & $(-6.6, -2.1, 18.3)$ & $(-4.4, -4.2, 11.3)$ \\
  corners a & $(-3.6, -2.4, 16.9)$ & \\
  corners c & $(-5.1, -0.6, 15.6)$ & \\
  corners b & $(6.1, -3.4, 10.7)$ & $(-1.0, -4.4, 9.6)$ \\
  corners d & $(3.2, -2.8, 10.0)$ & \\
\end{tabular}
\end{center}

Runs b and d drift smoothly, by $1$~cm every $5$ frames. They use $10$ to $12$ features, with no
outlier that matters, and report a position uncertainty of about $1$~cm. This is the
mechanism of the simulation. Everything that follows is a consequence of it. The features added
during the drift do not agree with those from the first part of the map. When the camera later
faces a door, old features are predicted in view behind the wall, and most of the searched
features fail. The map then freezes (section~\ref{sec:global}), and tracking runs out of features. The relocalizations then
jump between the two parts of the map, by about $45$~cm from one frame to the next.

\paragraph{The information is in the images.}
A bundle adjustment relinearizes every observation at the final estimate. Here it is run over the
$140$ poses and the new features, with the matches each run used, and with the pose at frame $860$
and the older features fixed. Fitted to the matches of run b alone, and to those of the default run alone, which are two
independent sets, the camera ends up $3.7$~cm apart at frame $1000$, where the two filters were
$14$~cm apart. The median reprojection errors are $0.5$ and $0.2$~px. The motion is in the images,
and the filter loses it, because it cannot revisit an update once it has been made.

\paragraph{Conclusion.}
This is the result of Strasdat et al.~\cite{strasdat2012}, which ended filtering as the method of
choice for visual SLAM. For the same computation, keyframes and bundle adjustment beat a filter,
because they can relinearize, and because they can afford many more points. Within a filter, the
remaining option is to add sensors: a gyroscope would measure the rotation, and the young features
could no longer confuse it with the translation. With a single camera, the next step is not a
better filter, but tracking against a map maintained by bundle adjustment, as in PTAM~\cite{klein2007}
and its successors.

\section{Correspondence with the code}

\begin{center}
\small
\begin{tabular}{ll}
  frame loop, parameters & \texttt{main.cpp} \\
  $R$, $\exp$, $[\cdot]_\times$ & \texttt{EKF::orientation}, \texttt{rotation}, \texttt{skew} in \texttt{camera.h} \\
  $f$, $F$, $W$, $Q$, $J_r$ & \texttt{MotionModel} in \texttt{ekf.h}, \texttt{rightJacobian} \\
  $\sigma$, update thresholds & \texttt{MeasurementModel} in \texttt{ekf.h} \\
  prediction & \texttt{EKF::predict} \\
  $h$, $H$, $S$ & \texttt{EKF::observe} \\
  update, 1-point RANSAC & \texttt{EKF::update}, \texttt{EKF::correct} \\
  state augmentation, removal & \texttt{EKF::addFeature}, \texttt{EKF::removeFeatures} \\
  relocalization of the camera & \texttt{EKF::resetCamera} \\
  inverse depth & \texttt{inverse\_depth::toCamera}, \texttt{fromPixel}, \texttt{fromPoint}, \texttt{toPoint} \\
  projection & \texttt{CameraModel::projectCameraPoint} \\
  initialization & \texttt{initializeFromPattern} \\
  patch, descriptors of a feature & \texttt{Feature} in \texttt{feature.h} \\
  seeds & \texttt{FlowTracker} in \texttt{tracking.h} \\
  patch search, returning features & \texttt{findPatch}, \texttt{matchFeatures} \\
  subpixel alignment, $\Sigma_a$ & \texttt{alignPatch} in \texttt{tracking.cpp} \\
  noise of each match & \texttt{matchNoise} in \texttt{main.cpp}, \texttt{EKF::setNoises} \\
  adding, removing features & \texttt{Mapper::addFeatures}, \texttt{Mapper::removeUnreliable} \\
  descriptors & \texttt{describe}, \texttt{Mapper::recordAppearance} \\
  keyframes & \texttt{Keyframes} in \texttt{relocalization.h} \\
  global matching & \texttt{matchGlobally} in \texttt{relocalization.h} \\
  input, output & \texttt{FrameSource}, \texttt{StateStream} \\
  window & \texttt{Display}, \texttt{draw} in \texttt{display.h} \\
\end{tabular}
\end{center}

\begin{thebibliography}{99}

\bibitem{baker2004}
S.~Baker and I.~Matthews.
Lucas-Kanade 20 years on: a unifying framework.
\emph{International Journal of Computer Vision}, 56(3):221--255, 2004.

\bibitem{civera2008}
J.~Civera, A.~J.~Davison and J.~M.~M.~Montiel.
Inverse depth parametrization for monocular SLAM.
\emph{IEEE Transactions on Robotics}, 24(5):932--945, 2008.

\bibitem{civera2010}
J.~Civera, O.~G.~Grasa, A.~J.~Davison and J.~M.~M.~Montiel.
1-point RANSAC for extended Kalman filtering: application to real-time structure from motion and
visual odometry.
\emph{Journal of Field Robotics}, 27(5):609--631, 2010.

\bibitem{davison2007}
A.~J.~Davison, I.~D.~Reid, N.~D.~Molton and O.~Stasse.
MonoSLAM: real-time single camera SLAM.
\emph{IEEE Transactions on Pattern Analysis and Machine Intelligence}, 29(6):1052--1067, 2007.

\bibitem{klein2007}
G.~Klein and D.~Murray.
Parallel tracking and mapping for small AR workspaces.
In \emph{Proceedings of the IEEE and ACM International Symposium on Mixed and Augmented Reality},
pages 225--234, 2007.

\bibitem{klein2008}
G.~Klein and D.~Murray.
Improving the agility of keyframe-based SLAM.
In \emph{Proceedings of the European Conference on Computer Vision}, pages 802--815, 2008.

\bibitem{lefferts1982}
E.~J.~Lefferts, F.~L.~Markley and M.~D.~Shuster.
Kalman filtering for spacecraft attitude estimation.
\emph{Journal of Guidance, Control, and Dynamics}, 5(5):417--429, 1982.

\bibitem{lucas1981}
B.~D.~Lucas and T.~Kanade.
An iterative image registration technique with an application to stereo vision.
In \emph{Proceedings of the 7th International Joint Conference on Artificial Intelligence}, pages
674--679, 1981.

\bibitem{rublee2011}
E.~Rublee, V.~Rabaud, K.~Konolige and G.~Bradski.
ORB: an efficient alternative to SIFT or SURF.
In \emph{Proceedings of the IEEE International Conference on Computer Vision}, pages 2564--2571,
2011.

\bibitem{shi1994}
J.~Shi and C.~Tomasi.
Good features to track.
In \emph{Proceedings of the IEEE Conference on Computer Vision and Pattern Recognition}, pages
593--600, 1994.

\bibitem{strasdat2012}
H.~Strasdat, J.~M.~M.~Montiel and A.~J.~Davison.
Visual SLAM: why filter?
\emph{Image and Vision Computing}, 30(2):65--77, 2012.

\bibitem{williams2011}
B.~Williams, G.~Klein and I.~Reid.
Automatic relocalization and loop closing for real-time monocular SLAM.
\emph{IEEE Transactions on Pattern Analysis and Machine Intelligence}, 33(9):1699--1712, 2011.

\end{thebibliography}

\end{document}
